\documentclass[rnd]{mas_proposal}
% \documentclass[thesis]{mas_proposal}

\usepackage[utf8]{inputenc}
\usepackage{amsmath}
\usepackage{amsfonts}
\usepackage{amssymb}
\usepackage{graphicx}

\title{Flow-Based Label Propagation for UAV Video Segmentation}
\author{Sunesh Praveen Raja Sundarasami}
\supervisors{Prof. Dr. Sebastian Houben}
\date{May 2026}

\begin{document}

\maketitle

\pagestyle{plain}

\section{Introduction}

\subsection{Topic of This DLRV project}

Semantic segmentation of UAV video is expensive because every frame requires
pixel-accurate labels when annotated manually. This project implements a
flow-based label propagation pipeline for a single UAV video from the
Ruralscapes dataset using SEA-RAFT~\cite{searaft2024} as the optical flow
backbone. Keyframe segmentation masks are warped to surrounding frames along
the estimated flow field and evaluated against dense ground truth using mIoU,
producing a clear picture of how label quality degrades with increasing
distance from the keyframe.

\subsection{Relevance of This DLRV project}

In a 30\,fps recording, one minute of video contains 1,800 frames. If the
current workflow labels one frame per second, that is already 60 manual
annotations per minute. A working propagation pipeline can extend those labels
to cover many more frames with acceptable quality, directly reducing annotation
cost for aerial video applications such as agricultural monitoring and
infrastructure inspection.

\section{Related Work}

\subsection{Survey of Related Work}

Flow-based label propagation has been studied extensively in the context of
video segmentation. The most directly related prior work is
SegProp~\cite{segprop2020}, which propagates sparse keyframe annotations to
unlabelled aerial video frames using optical flow, demonstrating that highly
redundant consecutive frames make this approach particularly effective for UAV
footage. The optical flow estimation in this project relies on
SEA-RAFT~\cite{searaft2024}, which extends the recurrent architecture of
RAFT~\cite{raft2020} with a simplified encoder and attention module, yielding
strong accuracy on standard benchmarks. The classical foundations of
flow-based warping and the forward-backward consistency check used here for
occlusion detection are rooted in the variational framework of Brox et
al.~\cite{brox2004}. Building on propagation as a core primitive, Zhu et
al.~\cite{zhu2019} propose a warp-and-refine paradigm in which
flow-propagated labels are subsequently corrected by a lightweight network,
providing motivation for the error analysis central to this project.

\subsection{Limitation and Deficits in the State of the Art}

Most propagation pipelines assume smooth, near-planar motion. UAV footage
violates this when the drone changes altitude or flies over scenes with depth
variation, introducing parallax that standard flow warping cannot fully
correct~\cite{segprop2020,zhu2019}. Furthermore, propagation quality is rarely
measured as a function of frame distance, leaving practitioners with little
guidance on how far from a keyframe the propagated labels remain
useful~\cite{raft2020,searaft2024}.

\section{Problem Statement}

This project builds a flow-based label propagation pipeline for a single
densely annotated UAV video and characterises how propagation quality degrades
with increasing distance from the keyframe. SEA-RAFT~\cite{searaft2024}
estimates optical flow from each keyframe to surrounding frames, the keyframe
mask is warped along the flow field using nearest-neighbour interpolation, and
occluded pixels are detected and excluded via a forward-backward consistency
check~\cite{brox2004}. The warped masks are evaluated against ground truth
using mIoU, and scores are plotted as a function of propagation distance to
directly characterise pipeline reliability.

\section{Project Plan}

\subsection{Work Packages}

\begin{enumerate}

    \item[WP1] \textbf{Literature Study and Setup.} Review the core literature
    on optical flow and flow-based label propagation, identify key design
    decisions from SegProp~\cite{segprop2020}, install all dependencies, and
    confirm that SEA-RAFT produces valid flow on at least one Ruralscapes frame
    pair.

    \item[WP2] \textbf{Propagation Pipeline.} Implement the end-to-end
    pipeline: run SEA-RAFT between keyframes and target frames, warp keyframe
    masks using nearest-neighbour interpolation, exclude occluded pixels via
    the forward-backward consistency check, and compute mIoU over valid pixels
    against ground truth. Validate on a short clip by confirming mIoU decreases
    with increasing frame distance.

    \item[WP3] \textbf{Evaluation, Analysis, and Report.} Scale the pipeline
    to the full video, record per-pair mIoU within a fixed propagation window,
    and produce a mIoU-vs-distance decay curve and a difficulty heatmap over
    the video timeline. Select a small set of failure cases and visualise them
    alongside ground-truth masks. Write the report continuously throughout the
    project.

\end{enumerate}

\subsection{Milestones}

\begin{enumerate}
    \item[M1] Literature reviewed; SEA-RAFT producing valid flow on one
    Ruralscapes frame pair.
    \item[M2] End-to-end pipeline validated on a short clip with correct mIoU
    evaluation.
    \item[M3] Full-video evaluation complete; heatmap, distance plot, failure
    cases, and final report submitted.
\end{enumerate}

\subsection{Deliverables}

\subsubsection*{Minimum Viable}
\begin{itemize}
    \item A working propagation pipeline with mIoU evaluation on one
    Ruralscapes video.
    \item A mIoU-vs-distance plot and at least two annotated failure-case
    visualisations.
\end{itemize}

\subsubsection*{Expected}
\begin{itemize}
    \item A difficulty heatmap over the full video with a written
    interpretation of failure patterns.
    \item A per-class IoU breakdown identifying which categories propagate
    most and least reliably.
\end{itemize}

\subsubsection*{Desired}
\begin{itemize}
    \item A clean, documented implementation reusable on other UAV videos.
    \item Practical guidelines for reliable propagation distance under
    different motion conditions.
\end{itemize}



\section{Potential Extension}

If the core pipeline is completed ahead of schedule, a frame selection model
may be pursued as an optional extension. SegProp~\cite{segprop2020} selects
keyframes at fixed intervals, ignoring content. The per-pair mIoU scores from
WP3 expose which intervals are hardest to propagate, providing a natural
signal for smarter selection. The extension would build a lightweight model
that automatically recommends which $k$ frames are most valuable to annotate
manually, prioritising frames at the boundary of high-difficulty intervals.
The goal is to show that content-aware keyframe selection yields higher overall
segmentation quality than the fixed-interval baseline of
SegProp~\cite{segprop2020} under the same annotation budget.

\nocite{*}

\bibliographystyle{unsrtnat}
\bibliography{bibliography.bib}

\end{document}